\subsection{Parametrické vyjádření přímky v prostoru}
\label{subsec:Parametric_line_in_space}

V předchozí části jsme si připomněli, že přímku můžeme určit jedním jejím
bodem a směrem. V prostoru budeme tento směr, stejně jako v rovině,
popisovat pomocí nenulového směrového vektoru.

Uvažujme tedy přímku $p$, která prochází bodem
\[
    A[x_A;y_A;z_A]
\]
a má směrový vektor
\[
    \vec{u}=(u_x;u_y;u_z).
\]

Zvolme nyní libovolný bod $X[x;y;z]$ této přímky. Protože body $A$ a $X$
leží na stejné přímce, je vektor
\[
    \overrightarrow{AX}
\]
rovnoběžný se směrovým vektorem $\vec{u}$. Musí tedy existovat reálné číslo
$t$ takové, že
\[
    \overrightarrow{AX}=t\vec{u}.
\]

Zároveň víme, že
\[
    \overrightarrow{AX}=X-A.
\]
Po dosazení proto dostáváme
\[
    X-A=t\vec{u},
\]
a tedy
\begin{equation}
    X=A+t\vec{u}.
    \label{eq:Parametric_line_in_space_vector}
\end{equation}

Tento vztah má jednoduchý geometrický význam. Začneme v pevném bodě $A$
a přičítáme k němu různé reálné násobky směrového vektoru $\vec{u}$.
Tím získáváme postupně všechny body přímky.

Po rozepsání vztahu~\eqref{eq:Parametric_line_in_space_vector} do
jednotlivých souřadnic dostáváme
\[
    [x;y;z]
    =
    [x_A;y_A;z_A]
    +
    t(u_x;u_y;u_z),
\]
tedy
\[
\begin{aligned}
    p: x&=x_A+t u_x,\\
    y&=y_A+t u_y,\\
    z&=z_A+t u_z; t\in\mathbb{R}.
\end{aligned}
\]

\begin{definition}
    \label{def:Parametric_line_in_space}
    Nechť $A[x_A;y_A;z_A]$ je bod v prostoru a
    $\vec{u}=(u_x;u_y;u_z)$ je nenulový vektor.
    Přímka $p$ procházející bodem $A$ a mající směrový vektor $\vec{u}$
    je množina všech bodů $X[x;y;z]$, které lze vyjádřit ve tvaru
    \[
        X=A+t\vec{u},\qquad t\in\mathbb{R}.
    \]
    V souřadnicích tedy
    \begin{equation}
        \begin{aligned}
            p: x&=x_A+t u_x,\\
            y&=y_A+t u_y,\\
            z&=z_A+t u_z; t\in\mathbb{R}.
        \end{aligned}
        \label{eq:Parametric_line_in_space_coordinates}
    \end{equation}

    Číslo $t$ nazýváme \textbf{parametrem} a vektor $\vec{u}$
    \textbf{směrovým vektorem} přímky.
\end{definition}

Podmínka $\vec{u}\neq\vec{0}$ je nezbytná. Kdyby byl směrový vektor nulový,
potom by pro každé $t\in\mathbb{R}$ platilo
\[
    X=A+t\vec{0}=A.
\]
Místo celé přímky bychom tak získali pouze jediný bod.

Geometrický význam parametru $t$ můžeme vidět na
obrázku~\ref{fig:Parametric_line_in_space}. Pro $t=0$ dostáváme přímo bod
$A$. Pro kladné hodnoty $t$ se z bodu $A$ pohybujeme ve směru vektoru
$\vec{u}$ a pro záporné hodnoty v opačném směru.

\begin{figure}[H]
    \centering
    \begin{tikzpicture}[scale=0.95]

        % Kladné části os
        \draw[->] (0,0,0) -- (5,0,0) node[right] {$x$};
        \draw[->] (0,0,0) -- (0,4.5,0) node[above] {$y$};
        \draw[->] (0,0,0) -- (0,0,4) node[below left] {$z$};

        % Záporné části os
        \draw[->, dashed] (0,0,0) -- (-2,0,0);
        \draw[->, dashed] (0,0,0) -- (0,-2,0);
        \draw[->, dashed] (0,0,0) -- (0,0,-2);

        \node[above left] at (0,0,0) {$O$};

        % Body ležící na přímce
        \coordinate (P1) at (-0.6,-0.2,0);
        \coordinate (Xm) at (0.6,0.7,0.75);
        \coordinate (A)  at (1.8,1.6,1.5);
        \coordinate (Xp) at (3,2.5,2.25);
        \coordinate (P2) at (4.2,3.4,3);
        \coordinate (P0) at (-1.5,-0.87,-0.56);

        % Přímka
        \draw[thick] (P1) -- (P2)
            node[pos=0.96, above] {$p$};
        \draw[thick, dashed] (P0) -- (P1);

        % Směrový vektor
        \draw[vector] (A) -- (Xp)
            node[midway, sloped, above] {$\vec{u}$};

        % Body
        \fill (A) circle (2pt)
            node[below right] {$A=X_0$};

        \fill (Xm) circle (1.7pt)
            node[above left] {$X_{-1}$};

        \fill (Xp) circle (1.7pt)
            node[below right] {$X_1$};

    \end{tikzpicture}
    \caption{Změnou hodnoty parametru $t$ získáváme jednotlivé body
    přímky $p$.}
    \label{fig:Parametric_line_in_space}
\end{figure}

Ze vztahu
\[
    \overrightarrow{AX}=t\vec{u}
\]
navíc vidíme, že absolutní hodnota parametru souvisí se vzdáleností bodu
$X$ od bodu $A$. Platí totiž
\[
    |\overrightarrow{AX}|
    =
    |t\vec{u}|
    =
    |t|\cdot|\vec{u}|.
\]
Například pro $t=2$ se tedy z bodu $A$ posuneme ve směru vektoru
$\vec{u}$ o dvojnásobek jeho velikosti, zatímco pro $t=-2$ o stejnou
vzdálenost v opačném směru.

\begin{example}
    \label{exa:Parametric_line_in_space_point_direction}
    Určete parametrické vyjádření přímky $p$, která prochází bodem
    \[
        A[1;-2;3]
    \]
    a má směrový vektor
    \[
        \vec{u}=(2;1;-1).
    \]

    Použijeme vztah
    \[
        X=A+t\vec{u}.
    \]
    Po dosazení získáme
    \begin{align*}
        X
        &=[1;-2;3]+t(2;1;-1)\\
        &=[1+2t;-2+t;3-t].
    \end{align*}

    Parametrické vyjádření přímky je tedy
    \[
        \begin{aligned}
            p: x&=1+2t,\\
            y&=-2+t,\\
            z&=3-t; t\in\mathbb{R},
        \end{aligned}
    \]
    nebo také ve vektorovém tvaru
    \[
        X=\left\{[1;-2;3]+t(2;1;-1); t\in\mathbb{R}\right\}.
    \]
\end{example}

Z parametrického vyjádření můžeme naopak snadno určit jeden bod přímky
a její směrový vektor. Konstantní členy odpovídají souřadnicím bodu, který
získáme pro $t=0$, zatímco koeficienty u parametru tvoří souřadnice
směrového vektoru.

\begin{example}
    \label{exa:Reading_line_in_space}
    Je dána přímka
    \[
        \begin{aligned}
            q: x&=-3+4t,\\
            y&=2-2t,\\
            z&=1+3t; t\in\mathbb{R}.
        \end{aligned}
    \]
    Určete jeden bod přímky $q$ a její směrový vektor.

    Pro $t=0$ dostáváme
    \[
        x=-3 + 4\cdot 0 = -3,\qquad y=2 - 2\cdot 0 = 2,\qquad z=1 + 3\cdot 0 = 1,
    \]
    a tedy například bod
    \[
        B[-3;2;1].
    \]

    Souřadnice směrového vektoru jsou dány koeficienty u parametru $t$:
    \[
        \vec{v}=(4;-2;3).
    \]

    Přímka $q$ tedy prochází bodem $B[-3;2;1]$ a má směrový vektor
    $\vec{v}=(4;-2;3)$.
\end{example}

Je důležité si uvědomit, že parametrické vyjádření přímky není určeno
jednoznačně. Bod $A$ v zápisu
\[
    X=A+t\vec{u}
\]
může být kterýkoli bod přímky a směrový vektor můžeme nahradit libovolným
jeho nenulovým reálným násobkem.

Je-li tedy $\vec{u}$ směrovým vektorem přímky, potom je jejím směrovým
vektorem také
\[
    k\vec{u},
    \qquad
    k\in\mathbb{R}\setminus\{0\}.
\]

Například přímku
\[
    \begin{aligned}
        p: x&=1+2t,\\
        y&=-2+t,\\
        z&=3-t; t\in\mathbb{R}
    \end{aligned}
\]
můžeme zapsat také pomocí bodu odpovídajícího hodnotě $t=2$:
\[
    B[5;0;1].
\]
Jako nový směrový vektor můžeme zvolit například
\[
    -2\vec{u}=(-4;-2;2).
\]
Dostáváme proto jiné parametrické vyjádření
\[
    \begin{aligned}
        p: x&=5-4s,\\
        y&=-2s,\\
        z&=1+2s; s\in\mathbb{R},
    \end{aligned}
\]
které však popisuje tutéž přímku.

Pomocí parametru můžeme také jednoduše určovat další body přímky. Stačí
za něj dosadit libovolné reálné číslo.

\begin{example}
    \label{exa:Points_on_parametric_line_space}
    Určete body přímky
    \[
        \begin{aligned}
            p: x&=1+2t,\\
            y&=-2+t,\\
            z&=3-t; t\in\mathbb{R}
        \end{aligned}
    \]
    odpovídající hodnotám $t=-1$, $t=0$ a $t=2$.

    Pro $t=-1$ dostáváme
    \[
        X_{-1}
        =
        [1+2(-1);-2+1(-1);3-1(-1)]
        =
        [-1;-3;4].
    \]

    Pro $t=0$ dostáváme
    \[
        X_0 
        = 
        [1+2(0);-2+1(0);3-1(0)] 
        =
        [1;-2;3].
    \]

    Pro $t=2$ dostáváme
    \[
        X_2
        =
        [1+2(2);-2+1(2);3-1(2)]
        =
        [5;0;1].
    \]
\end{example}

Opačnou úlohou je rozhodnout, zda zadaný bod na přímce leží. Bod
$P[x_P;y_P;z_P]$ leží na přímce právě tehdy, když existuje jedna hodnota
parametru $t$, pro kterou současně platí všechny tři parametrické rovnice.

\begin{example}
    \label{exa:Point_membership_line_space}
    Rozhodněte, zda body
    \[
        B[5;0;1]
        \qquad\text{a}\qquad
        C[3;0;2]
    \]
    leží na přímce
    \[
        \begin{aligned}
            p: x&=1+2t,\\
            y&=-2+t,\\
            z&=3-t; t\in\mathbb{R}.
        \end{aligned}
    \]

    Nejprve dosadíme souřadnice bodu $B[5;0;1]$:
    \begin{align*}
        5&=1+2t &&\Rightarrow t=2,\\
        0&=-2+t &&\Rightarrow t=2,\\
        1&=3-t &&\Rightarrow t=2.
    \end{align*}

    Ze všech tří rovnic vyšla stejná hodnota parametru $t=2$, a proto
    bod $B$ na přímce leží:
    \[
        B\in p.
    \]

    Nyní dosadíme souřadnice bodu $C[3;0;2]$:
    \begin{align*}
        3&=1+2t &&\Rightarrow t=1,\\
        0&=-2+t &&\Rightarrow t=2,\\
        2&=3-t &&\Rightarrow t=1.
    \end{align*}

    Neexistuje jediná hodnota parametru, která by současně vyhovovala
    všem třem rovnicím. Proto bod $C$ na přímce neleží:
    \[
        C\notin p.
    \]
\end{example}

Stejnou podmínku můžeme vyjádřit také pomocí vektorů. Je-li přímka určena
bodem $A$ a směrovým vektorem $\vec{u}$, potom bod $P$ leží na této
přímce právě tehdy, když je vektor
\[
    \overrightarrow{AP}
\]
reálným násobkem směrového vektoru $\vec{u}$, tedy když existuje
$t\in\mathbb{R}$ takové, že
\[
    \overrightarrow{AP}=t\vec{u}.
\]

Přímka nemusí být vždy zadána bodem a směrovým vektorem. Velmi často
známe dva různé body
\[
    A[x_A;y_A;z_A],
    \qquad
    B[x_B;y_B;z_B],
\]
kterými má přímka procházet.

Protože dvěma různými body prochází právě jedna přímka, můžeme její
směrový vektor získat přímo jako
\[
    \overrightarrow{AB}
    =
    (x_B-x_A;y_B-y_A;z_B-z_A).
\]

Dosadíme-li tento vektor do vztahu $X=A+t\vec{u}$, dostáváme
\[
    X=A+t\overrightarrow{AB}.
\]
Po rozepsání do souřadnic tedy
\[
    \begin{aligned}
        p: x&=x_A+t(x_B-x_A),\\
        y&=y_A+t(y_B-y_A),\\
        z&=z_A+t(z_B-z_A); t\in\mathbb{R}.
    \end{aligned}
\]

\begin{example}
    \label{exa:Parametric_line_in_space_two_points}
    Určete parametrické vyjádření přímky $r$, která prochází body
    \[
        K[-1;2;4]
        \qquad\text{a}\qquad
        L[3;-2;1].
    \]

    Nejprve určíme směrový vektor přímky:
    \[
        \overrightarrow{KL}
        =
        (3-(-1);-2-2;1-4)
        =
        (4;-4;-3).
    \]

    Použijeme například bod $K[-1;2;4]$ a směrový vektor
    $\overrightarrow{KL}=(4;-4;-3)$. Dostáváme
    \[
        \begin{aligned}
            r: x&=-1+4t,\\
            y&=2-4t,\\
            z&=4-3t; t\in\mathbb{R}.
        \end{aligned}
    \]

    Pro $t=0$ získáme bod $K$ a pro $t=1$ bod
    \[
        [-1+4;2-4;4-3]
        =
        [3;-2;1]
        =
        L.
    \]
    Přímka tedy skutečně prochází oběma zadanými body.
\end{example}

Při určování přímky dvěma body bychom mohli stejně dobře použít bod $L$
a opačný vektor
\[
    \overrightarrow{LK}
    =
    -\overrightarrow{KL}.
\]
Výsledné parametrické vyjádření by vypadalo jinak, ale popisovalo by
stejnou přímku.